\documentclass[11pt,a4paper]{article}

\usepackage[margin=1in]{geometry}
\usepackage{amsmath,amssymb,amsthm}
\usepackage{hyperref}
\usepackage{booktabs}
\usepackage{listings}
\usepackage{xcolor}
\usepackage{tikz}
\usepackage{enumitem}

\hypersetup{colorlinks=true,linkcolor=blue,citecolor=blue,urlcolor=blue}

\newtheorem{theorem}{Theorem}
\newtheorem{lemma}[theorem]{Lemma}
\newtheorem{corollary}[theorem]{Corollary}
\newtheorem{definition}[theorem]{Definition}
\newtheorem{proposition}[theorem]{Proposition}

\title{ZAP-WS: Per-Stream Independent Flow Control with\\
Tail-Latency Guarantees and Post-Quantum Confidentiality}

\author{ZAP Protocol Authors\\
\texttt{zap-proto.io}}

\date{2026}

\begin{document}
\maketitle

\begin{abstract}
We present ZAP-WS, a streaming and publish/subscribe layer over the ZAP transport
protocol that provides multiple independent logical streams over a single ZAP
connection. Unlike WebSocket (RFC 6455) and HTTP/2 (RFC 9113), which both
multiplex streams over a single TCP byte-stream and therefore inherit
TCP-level head-of-line (HoL) blocking on packet loss, ZAP-WS gives each
stream its own sequence number space, frame-level optional forward error
correction (FEC), and per-stream credit-based backpressure. The result is a
tail-latency profile in which the arrival latency of stream $s_i$ is
statistically independent of frames in stream $s_j$ for $j \ne i$. ZAP-WS
matches the head-of-line blocking elimination of HTTP/3 (RFC 9114) over
QUIC, while retaining mutual authentication and post-quantum
confidentiality from the underlying ZAP handshake. We formalize the
tail-latency claim and prove it under standard packet loss models. We
discuss when ZAP-WS is materially better than HTTP/3 (post-quantum
deployments, embedded environments without QUIC, scenarios requiring
deterministic FEC instead of retransmission).
\end{abstract}

\section{Introduction}

Modern application traffic is dominated by long-lived streaming patterns:
chat, market data, telemetry, live updates, agent-to-agent coordination,
log shipping, and progressive AI inference. The transport options
historically available to application developers are constrained:

\begin{itemize}
\item \textbf{Server-Sent Events (SSE)} -- unidirectional (server $\to$
client), text-only, single stream per HTTP connection.
\item \textbf{WebSocket (RFC 6455)} \cite{rfc6455} -- full-duplex, but a
single ordered byte stream per TCP connection. Application-level
multiplexing requires explicit framing and forfeits any independence
guarantee.
\item \textbf{HTTP/2 (RFC 9113)} \cite{rfc9113} -- multi-stream
multiplexing over a single TCP connection. The streams share TLS and TCP
state; a single packet loss blocks delivery of \emph{all} streams until
retransmission, even though the streams are logically independent.
\item \textbf{HTTP/3 (RFC 9114)} \cite{rfc9114} over QUIC (RFC 9000)
\cite{rfc9000} -- multi-stream with per-stream loss recovery, eliminating
TCP-level HoL blocking. Adoption lags in embedded, kernel-bypass,
broker-to-broker, and trusted-network environments. Post-quantum profiles
are still being standardized.
\end{itemize}

ZAP-WS occupies the same design space as HTTP/3 streams but differs in
three respects: (i) the underlying ZAP handshake is post-quantum by
default (ML-KEM-768 hybrid with X25519, ML-DSA-65 with Ed25519); (ii) the
per-stream loss recovery mechanism may use Reed--Solomon FEC rather than
retransmission, eliminating one round-trip on the recovery path;
(iii) the API is exposed as a $\texttt{net.Conn}$-shaped object per
stream, so existing WebSocket, $\texttt{io.Reader}/\texttt{io.Writer}$,
and SCTP-style code ports mechanically.

\paragraph{Thesis.} ZAP-WS provides per-stream independent flow control,
eliminating the head-of-line blocking that affects WebSocket and HTTP/2
streams sharing a TCP connection, while retaining mutual authentication
and post-quantum confidentiality from the underlying ZAP transport.

\paragraph{Contributions.}
\begin{enumerate}
\item A frame format and stream multiplexing scheme that gives each
logical stream its own sequence space, window, and (optional) FEC group.
\item A formal tail-latency independence theorem (Theorem 1) and a
throughput composition theorem (Theorem 2), proved under a Bernoulli
packet-loss model.
\item A practical comparison with WebSocket, HTTP/2, HTTP/3, MPTCP
\cite{rfc8684}, and SCTP \cite{rfc4960}.
\item A reference implementation sketch and guidance on when ZAP-WS is
materially better than HTTP/3.
\end{enumerate}

\section{Head-of-Line Blocking Analysis}
\label{sec:hol}

\subsection{TCP byte-stream model}

A TCP connection delivers bytes in order. Internally, TCP buffers
out-of-order segments at the receiver and waits for missing segments to
be retransmitted before delivering subsequent in-order data to the
application. Concretely, if segments
$S_1, S_2, \ldots, S_n$ are sent and $S_k$ is lost, segments
$S_{k+1}, \ldots, S_n$ are buffered at the receiver until $S_k$ is
retransmitted and arrives. The application sees a stall of one
retransmission timeout (RTO) or fast-retransmit RTT, regardless of
whether the bytes in $S_{k+1}, \ldots, S_n$ are part of the same logical
unit of work as $S_k$.

\subsection{HTTP/2 over TCP}

HTTP/2 multiplexes \emph{streams} -- bidirectional sequences of frames --
over a single TCP connection \cite{rfc9113}. Frames from different
streams are interleaved on the wire. A lost segment carrying frames of
stream 7 will, by the TCP semantics above, block delivery of any stream's
frames sent in subsequent segments. Stream 9's frame may already have
arrived at the kernel but the application cannot read it.

\subsection{HTTP/3 over QUIC}

QUIC implements stream multiplexing inside the cryptographic transport
itself \cite{rfc9000}. Each stream has independent flow control and
independent loss detection. A lost packet only delays the streams whose
frames it actually contained. HoL blocking at the transport layer is
eliminated.

\subsection{ZAP-WS}

ZAP-WS adopts the QUIC philosophy -- streams are first-class transport
objects, not application-layer multiplexes over a byte stream -- but
implements them on top of ZAP datagrams (which themselves run over UDP
or other packet substrates). Each ZAP-WS stream has:

\begin{enumerate}
\item An independent 64-bit sequence space.
\item An independent receive window and credit-based backpressure
($\texttt{WINDOW\_UPDATE}$ frames sent per stream, not per connection).
\item An optional FEC group: a $(k, m)$ Reed--Solomon code applied
across $k$ data frames producing $m$ parity frames, allowing the
receiver to reconstruct up to $m$ losses without round-trip
retransmission.
\end{enumerate}

\section{ZAP-WS Construction}
\label{sec:construction}

\subsection{Frame format}

Each ZAP-WS frame on the wire is:

\begin{verbatim}
  +------+--------+--------+----------+----------+--------+
  | Type | StrmID | StrmSeq| FECGrp   | FECIndex | Payld  |
  | 1B   | varint | varint | varint?  | varint?  | bytes  |
  +------+--------+--------+----------+----------+--------+
\end{verbatim}

Frame types: \texttt{DATA}, \texttt{WINDOW\_UPDATE},
\texttt{STREAM\_OPEN}, \texttt{STREAM\_CLOSE}, \texttt{FEC\_PARITY},
\texttt{PING}, \texttt{GOAWAY}. The \texttt{FECGrp} and
\texttt{FECIndex} fields are present only when the connection negotiated
FEC for the stream.

\subsection{Stream lifecycle}

A stream is opened by either peer with \texttt{STREAM\_OPEN}. Streams are
identified by 62-bit IDs. Bit 0 distinguishes initiator (client/server),
bit 1 distinguishes reliable/unreliable. After \texttt{STREAM\_CLOSE}
from both sides, the stream ID is retired but never reused on the same
connection.

\subsection{Per-stream flow control}

Each receiver advertises an initial window $W_0$ (typically 256~KiB) and
sends \texttt{WINDOW\_UPDATE} frames as the application drains data. A
sender that exhausts a stream's window blocks \emph{only that stream};
other streams are unaffected.

\subsection{Forward error correction}

When FEC is negotiated for a stream with parameters $(k, m)$, the sender
groups $k$ consecutive \texttt{DATA} frames into a FEC group, computes
$m$ \texttt{FEC\_PARITY} frames using a systematic Reed--Solomon code
\cite{reed1960}, and transmits the $k+m$ frames. The receiver can
reconstruct the original $k$ data frames after receiving any $k$ of the
$k+m$, tolerating up to $m$ losses without retransmission.

\paragraph{Overhead bound.} The bandwidth overhead of $(k, m)$ FEC is
exactly $m/k$. This is independent of which frames are lost.

\subsection{Cryptographic envelope}

ZAP-WS frames travel inside ZAP records that are AEAD-encrypted under a
session key established by the ZAP handshake \cite{zap-handshake}. The
handshake is hybrid post-quantum: ML-KEM-768 + X25519 for key
encapsulation and ML-DSA-65 + Ed25519 for authentication. Confidentiality
and authenticity therefore hold under both classical and quantum
adversaries (assuming ML-KEM and ML-DSA assumptions). HTTP/3 does not yet
mandate post-quantum profiles; ZAP-WS does, by construction.

\section{Comparison}
\label{sec:comparison}

\begin{table}[h]
\centering
\small
\begin{tabular}{lcccc}
\toprule
                       & WebSocket & HTTP/2 & HTTP/3 & ZAP-WS \\
\midrule
Multi-stream           & app-layer & native & native  & native \\
TCP HoL blocking       & yes       & yes    & n/a     & n/a \\
Per-stream window      & no        & yes    & yes     & yes \\
Loss recovery          & TCP rtx   & TCP rtx& QUIC rtx& rtx + FEC \\
Post-quantum mandatory & no        & no     & no      & yes \\
Embedded-friendly      & yes       & no     & limited & yes \\
\bottomrule
\end{tabular}
\caption{Streaming protocol comparison.}
\end{table}

\section{Latency Model}
\label{sec:latency}

Let a connection carry $N$ streams $s_1, \ldots, s_N$. Let
$L(s_i, f)$ denote the application-observed arrival latency of frame $f$
of stream $s_i$. Assume an i.i.d.\ Bernoulli packet-loss channel with
loss probability $p$ and one-way base delay $D$.

\begin{proposition}[WebSocket / HTTP/2 over TCP -- HoL]
For WebSocket and HTTP/2, the latency of stream $s_i$'s frame $f$
satisfies
\[
L(s_i, f) \ge D + \mathbb{1}[\text{any earlier byte in TCP order was
lost}] \cdot R,
\]
where $R$ is the retransmission delay (bounded below by one RTT, often
RTO). Hence $L(s_i, f)$ is dependent on traffic of all other streams
$s_j$, $j \ne i$.
\end{proposition}

\begin{proposition}[ZAP-WS / HTTP/3 -- no HoL]
For ZAP-WS, the latency of stream $s_i$'s frame $f$ satisfies
\[
L(s_i, f) \ge D + \mathbb{1}[\text{a packet carrying part of $f$ in
$s_i$ was lost}] \cdot R',
\]
where $R'$ is the per-stream recovery delay -- either an RTT for
retransmission, or $0$ if FEC parity within the stream's FEC group is
sufficient to reconstruct $f$. The indicator depends only on losses
affecting $s_i$.
\end{proposition}

These propositions are formalized and proved in
$\texttt{proofs/streaming-tail-latency/proof.tex}$ as Theorem 1.

\section{Use Cases}

\subsection{Market data fanout}

A market-data publisher serves $N$ symbols to many subscribers. With
WebSocket, a single packet drop in the symbol-A stream stalls symbol-B
quotes. With ZAP-WS, each symbol can be its own stream with its own FEC
group, sized to the symbol's update cadence and tolerable staleness.
Tail-latency p99.9 on symbol B is independent of symbol-A loss events.

\subsection{Agent-to-agent task streams}

In multi-agent systems, an orchestrator drives many agents in parallel.
Each agent's task plan, tool calls, and partial results are a logical
stream. Per-stream independence means a slow agent (long tool execution,
local disk wait) cannot block fast agents' frames.

\subsection{Telemetry / log shipping}

Distinct services shipping logs over a single ZAP connection get
independent backpressure. A burst of error-log volume from service X does
not delay metric frames from service Y.

\subsection{Real-time notifications}

Notification services often multiplex many subscriptions per user. With
ZAP-WS, a stalled subscription does not delay the rest of the user's
subscriptions.

\section{Implementation}

The reference ZAP-WS implementation exposes the following Go API
(adapted from the existing $\texttt{net.Conn}$ interface):

\begin{lstlisting}[language=Go, basicstyle=\small\ttfamily]
type Conn interface {
    OpenStream(opts StreamOpts) (Stream, error)
    AcceptStream() (Stream, error)
    Close() error
}

type Stream interface {
    net.Conn   // Read, Write, Close, deadlines
    StreamID() uint64
    SetFEC(k, m int) error
}
\end{lstlisting}

This signature is intentionally compatible with $\texttt{golang.org/x/net/quic.Stream}$
shape; existing code that writes to a QUIC stream ports without
substantive change. Bindings exist for C, Rust, Python, Java, JS,
OCaml, and the other targets in the $\texttt{zap-proto.io}$ family.

\section{When NOT to use ZAP-WS}

ZAP-WS is targeted at \emph{many concurrent logical streams}. For pure
bulk file transfer where the natural shape is a single ordered byte
stream, the ZAP record interface (no per-stream framing overhead) is
preferable. For scenarios where a single very large message is sent
periodically with no concurrent traffic, WebSocket or HTTP/2 are
adequate; the FEC and per-stream overhead of ZAP-WS would add bytes
without latency benefit.

\section{When ZAP-WS is materially better than HTTP/3}

\begin{enumerate}
\item \textbf{Post-quantum mandates.} HTTP/3 standardization is still
phasing in PQ key exchange. ZAP-WS is PQ-by-default today.
\item \textbf{FEC-driven tail control.} HTTP/3 recovers lost packets via
retransmission, costing at least one RTT. ZAP-WS optionally uses
Reed--Solomon FEC, recovering losses with zero added round trips at the
cost of $m/k$ bandwidth. For p99.9 latency in lossy networks
(satellite, mobile, congested datacenter east-west), this matters.
\item \textbf{Embedded / kernel-bypass.} QUIC implementations are large
and userspace-coupled; ZAP-WS has small implementations suitable for
microcontrollers and DPDK fast paths.
\item \textbf{Trusted-network broker meshes.} ZAP-WS provides
mutual-auth-by-default with PQ identities, which simplifies broker
mesh deployments compared to bolting mTLS onto HTTP/3.
\end{enumerate}

\section{Related Work}

WebSocket \cite{rfc6455}, HTTP/2 \cite{rfc9113}, HTTP/3 \cite{rfc9114},
QUIC \cite{rfc9000}, Reed--Solomon coding \cite{reed1960}, MPTCP
\cite{rfc8684}, SCTP \cite{rfc4960}, Google QUIC measurements
\cite{langley2017quic}. Head-of-line blocking has been measured
empirically by Wang et al.\ \cite{wang2014http2hol} and others.

\section{Conclusion}

ZAP-WS gives applications independent logical streams over a single
authenticated, post-quantum-encrypted connection. The independence is
formal: tail latency on one stream is statistically independent of
traffic on other streams (Theorem 1, proof in companion document).
Aggregate throughput composes additively up to link capacity (Theorem 2).
The construction is simple, ports existing WebSocket-style code without
API friction, and removes the head-of-line blocking that has been the
single biggest practical objection to multiplexing many real-time
streams over one connection.

\bibliographystyle{plain}
\bibliography{refs}

\end{document}
